\note{In each iteration, the generator outputs a candidate network trace, and the execution module runs the target and, when the reference is another protocol, runs the reference inside an extended Mahimahi emulator (which the authors modified to support time-varying one-way delay and dual paths). The scoring module then quantifies their relative performance gap as a \textbf{score}:\par \par \[
\text{score} = \frac{R - P}{\max(R, P)} \quad (\text{when higher metric is better}).
\]\par \par For example, if the reference achieves $R = 10$ Mbps and the target gets only $P = 2$ Mbps, the score is $(10 - 2)/10 = 0.8$. If a lower metric is better—such as flow completion time—we reverse the numerator to $P - R$. For nonnegative metrics with a positive denominator, dividing by the maximum normalizes the score between $-1$ and $1$, so a larger positive score always means a larger relative weakness in the target. We can also combine throughput $\tau$ and average packet delay $d$ into a single utility metric using a throughput weight $t_{\mathrm{coeff}} \in [0, 1]$:\par \par \[
U = t_{\mathrm{coeff}} \frac{\tau}{\tau_{\max}} + (1 - t_{\mathrm{coeff}}) \frac{d_{\min}}{d},
\]\par \par where $t_{\mathrm{coeff}} = 1$ cares only about throughput, $0$ cares only about delay, and $0.5$ weights the two normalized terms equally. Here, $\tau_{\max}$ and $d_{\min}$ are the throughput and delay normalization benchmarks.
\par\smallskip Source: AdvNet paper §3.5, Eq.(1)}
